\chapter{Experimental Design}
\label{chap:design}

The simulations draw count tables from populations with known MI. This
allows us to measure errors in the estimated MI difference and its
standard error, and to count false positives when the population MI
values are equal. The evaluation follows three questions: are the
variance approximations used by Expanded Welch accurate, does its
cutoff improve false-positive rates, and what happens to power when
false-positive rates are comparable? The design follows the simulation
framework of \citet{morris2019}, specifying the populations, quantities
being estimated, methods and performance measures.

\section{Constructing joint distributions with known MI}

The construction changes association while preserving the category
frequencies. For one population, draw \(X\) with category probabilities
\(a=(a_1,\ldots,a_k)\). With probability \(\lambda\), set \(Y=X\);
otherwise draw \(Y\) independently with the same probabilities \(a\).
Copying makes matching categories more likely than under independence.
Both branches give \(Y\) the same marginal probabilities, so changing
\(\lambda\) does not change the row or column totals.

An observation in a diagonal cell can arrive through either branch,
whereas an off-diagonal observation can arrive only through the
independent branch.
Combining these probabilities gives
\begin{equation}
 p_{ij}(\lambda)=(1-\lambda)a_i a_j
                  +\lambda a_i\mathbf{1}_{\{i=j\}},
 \qquad 0\leq\lambda<1,
 \label{eq:copying-construction}
\end{equation}
where the indicator is one when the two categories match and zero otherwise.
This construction is a special case of the mixture model described by
\citet{bocci2010}.
At \(\lambda=0\), the variables are independent and MI is zero.
Increasing \(\lambda\) increases MI, so the target MI determines the required
copying probability. Appendix~\ref{app:construction} gives the numerical
search and shows why the margins remain fixed.

\begin{samepage}
\medskip
\thesisminiheading{Constructing a \(2\times2\) joint distribution with positive MI.}

For \(P\), take \(a=(0.7,0.3)\) and a copying probability of
approximately 0.20347. Cell \((1,1)\) requires \(X=1\), which has
probability 0.7. Given \(X=1\), copying always gives \(Y=1\), whereas
an independent draw gives \(Y=1\) with probability 0.7. Thus
\begin{align*}
 p_{11}&\simeq0.7\{0.20347+(1-0.20347)(0.7)\}=0.53273,\\
 p_{12}&\simeq0.7(1-0.20347)(0.3)=0.16727.
\end{align*}
\end{samepage}
The second equation has no copying term because copying cannot produce
\(X=1,Y=2\). The remaining cells follow the same rule.
For \(Q\), use \(a=(0.8,0.2)\) and a copying probability of approximately
0.26091. These choices give \(I(P)=0.02\) and \(I(Q)=0.03\) nats:
\begin{equation}
 P\simeq\begin{pmatrix}0.53273&0.16727\\0.16727&0.13273\end{pmatrix},
 \qquad
 Q\simeq\begin{pmatrix}0.68175&0.11825\\0.11825&0.08175\end{pmatrix}.
 \label{eq:design-example-pair}
\end{equation}
Both joint distributions have positive MI, and the true difference used
to assess power has magnitude 0.01 nats.

\medskip
\thesisminiheading{Table sizes and margin regimes.}

The copying construction is used for square tables of sizes
\(2\times2,4\times4,8\times8\) and \(16\times16\), containing
4, 16, 64 and 256 cells. Each doubling of the number of categories
multiplies the number of cells by four, so at a fixed sample size the
observations are distributed over more possible cells.

Three margin regimes vary the category frequencies: uniform frequencies
in both populations, the same uneven frequencies in both, or different
uneven frequencies between them. For uneven margins, the first category
has probability \(\rho\), and the remaining mass is shared equally:
\begin{equation}
 m_k(\rho)=\left(\rho,\frac{1-\rho}{k-1},\ldots,
                   \frac{1-\rho}{k-1}\right).
 \label{eq:skew-margin}
\end{equation}
Each population uses the same vector for its row and column margins.
Table~\ref{tab:margin-profiles} specifies the vectors for \(P\) and \(Q\).
The copying probabilities for \(P\) and \(Q\) are chosen separately to
achieve the targets \(I(P)\) and \(I(Q)\).

\begin{table}[htbp]
\centering
\caption{Row and column margins for each population.}
\label{tab:margin-profiles}
\begin{tabular}{@{}lll@{}}
\toprule
Regime & \(P\): rows and columns & \(Q\): rows and columns\\
\midrule
Uniform & \((1/k,\ldots,1/k)\) & \((1/k,\ldots,1/k)\)\\
Same skew & \(m_k(0.8)\) & \(m_k(0.8)\)\\
Different skew & \(m_k(0.7)\) & \(m_k(0.8)\)\\
\bottomrule
\end{tabular}
\end{table}

Under equal MI, the uniform and same-skew regimes give identical joint
probability tables, \(P=Q\). In contrast, the different-skew regime gives
different tables with equal MI because the category frequencies differ
between populations. Comparing the regimes therefore tests false-positive
behaviour under both identical and distinct population tables.

With equal MI and equal sample sizes,
the same-skew samples have identical distributions, so their expected MI
estimation biases cancel in the difference. Different-skew populations can
retain unequal biases despite having equal MI. Under an MI difference or
unequal sample sizes, matching margins alone does not ensure that biases
cancel.

Increasing table size under the uneven-margin rule spreads the remaining
mass over more categories, changing both the number of cells and their
probabilities. Uniform versus same skew instead compares even
and uneven category frequencies at a fixed size; different skew allows the
two populations' frequencies to differ as well.



\section{Checking the variance approximations}

Expanded Welch chooses its degrees of freedom from the mean and sampling
variance of the estimated MI variance. For a fixed population \(P\),
\(\Vhat(P)\) estimates the pointwise-MI variance \(V(P)\), and
\(\Vhat(P)/n_P\) estimates that population's contribution to the variance
of the MI difference. Repeated samples also change \(\Vhat(P)\) itself.
Chapter~\ref{chap:expanded} approximates its two sampling moments by
\begin{equation}
 \E\{\Vhat(P)\}\approx V(P),\qquad
 \Var\{\Vhat(P)\}\approx\frac{\tau^2(P)}{n_P}.
 \label{eq:redesign-moment-approximations}
\end{equation}
Here \(\tau^2(P)\) is the variance of the observation-level sensitivity
derived in that chapter. The experiment compares each prediction with
the corresponding mean or variance measured across repeated samples.
It also checks the average sample-based prediction
\(\widehat\tau^2(P)/n_P\), because Expanded Welch estimates the sensitivity
from the observed count table.

These checks use common MI 0.02 nats across the four table sizes, three
margin regimes and sample sizes from 10 to 50,000. A separate
\(4\times4\) sweep varies the common positive MI from 0.0001 to 0.05
nats at \(n=100\), 1,000 and 10,000. Small positive MI approaches the
independence boundary, where the first-order variance becomes small and
quadratic sampling terms can matter. The complete grids appear in
Appendix~\ref{app:construction}.

To assess what moment errors do to the degrees of freedom, the experiment
also evaluates the moment-matching formula using the measured mean and
variance of \(\Vhat\). This gives a sampling-moment target for comparison
with the population first-order formula and the estimated component
degrees of freedom. Their mean-to-target ratio measures agreement,
while medians and percentile ranges describe sampling spread. Included
fractions identify any conditional summaries. The definitions and
algebra relating these quantities are in
Appendix~\ref{sec:moment-reference-details}. The near-independence
predictions are derived in Appendix~\ref{sec:independence-boundary}.

\section{False positives when MI is equal}

The main comparison uses Normal Wald, Expanded Welch and a sample-size
degrees-of-freedom rule. All three use the bias-corrected MI difference
and standard error from Chapter~\ref{chap:baselines}. Normal Wald uses
the standard normal distribution, while Expanded Welch uses the Student
\(t\) distribution with its MI-specific estimated degrees of freedom.
The sample-size rule asks whether the variance-sensitivity calculation
improves on assigning each component its sample size, as in
\citet{hutcheson1970}. It uses
\begin{equation}
 \widehat\nu_{\mathrm{size}}=
 \frac{\{\Vhat(P)/n_P+\Vhat(Q)/n_Q\}^{2}}
 {\{\Vhat(P)/n_P\}^{2}/n_P+\{\Vhat(Q)/n_Q\}^{2}/n_Q}.
 \label{eq:redesign-size-df}
\end{equation}
The component assignment is \(n_P\) and \(n_Q\); ordinary Welch's test
of sample means instead uses \(n_P-1\) and \(n_Q-1\)
\citep{welch1947}.

Each rule uses a two-sided 5\% nominal test. The broad null comparison
fixes \(I(P)=I(Q)=0.02\) nats and crosses the four table sizes, three
margin regimes and ten sample sizes, giving 120 settings. A rate above
5\% means too many false positives, while a rate below 5\% means
conservative rejection. Valid-result frequencies are reported alongside
these rates because an unusable calculation produces no rejection.

A further diagnostic uses a fixed Student \(t\) reference whose degrees
of freedom are fitted from independent null simulations. Comparing it
with Expanded Welch shows what happens when the sample-dependent degrees
of freedom are replaced by one value. The fixed reference appears in
the plots and tables alongside the three main rules. Its construction
and fitting procedure are given in
Appendix~\ref{sec:moment-reference-details}.

\section{Power when MI differs}

A lower false-positive rate can accompany lower power. The comparison
first uses nominal 5\% cutoffs, then adjusts them to match 5\% rejection
at a specified equal-MI population. Matching one null population does
not ensure the same false-positive rate elsewhere under equal MI.

The focused comparison uses the three populations in
Table~\ref{tab:redesign-focus-settings}. Settings A and B were selected
from earlier simulations because Expanded Welch brought false-positive
rates closer to 5\%; their power is evaluated using new independent
samples. Setting C provides a larger-sample comparison.

\begin{table}[!htbp]
\centering
\caption{Focus settings for the power comparison; common null MI is 0.02 nats.}
\label{tab:redesign-focus-settings}
\begin{tabular}{@{}llrl@{}}
\toprule Setting & Margins & \(n_P=n_Q\) & Size\\\midrule
A & Uniform & 20 & \(4\times4\)\\
B & Different skew & 100 & \(8\times8\)\\
C & Different skew & 10,000 & \(4\times4\)\\
\bottomrule
\end{tabular}
\end{table}

For each setting, an independent null simulation estimates the sampling
SD of the corrected MI difference. The power curve keeps \(I(P)=0.02\)
nats and increases \(I(Q)\) by zero to four of these null SDs, in steps
of half an SD. Scaling the differences this way makes the three curves
comparable relative to their sampling variation. The true signed
difference \(\Delta=I(P)-I(Q)\) is negative at every nonzero step.
Appendix~\ref{app:construction} gives the fixed SD values and resulting
population construction.

Cutoffs for each full curve are fitted using an independent batch with
\(I(P)=I(Q)=0.02\) nats. The same fitted cutoff is then applied
throughout that alternative curve. Separate null evaluation samples
check the achieved false-positive rate. Normal Wald and the fixed
moment reference give identical calibrated decisions because both order
the samples by \(|T|\). Expanded Welch and the sample-size rule can
order them differently because their degrees of freedom depend on the
sample. The fitted cutoffs evaluate the rules under known populations;
they are not an additional procedure for analysing arbitrary observed
tables.

A sensitivity check repeats calibration at two and four null SDs for
A--C. The lower null sets both MI values to 0.02; the upper null sets
both to the fixed alternative's larger MI, retaining margins and sample
sizes. Fresh upper-null fitting and validation samples allow both
calibrations to be compared on the same alternative pairs and checked
at both nulls. Appendix~\ref{sec:calibration-sensitivity-design} gives
the construction and uncertainty calculation.

The broader power comparison uses MI differences from 0 to 0.2 nats at
\(n=1{,}000\) and 10,000. Two further checks change the population
construction or the sample allocation. The construction check replaces
copying by a probability transfer within four cells, preserving the
margins and target MI. The allocation check keeps each population pair
and 2,000 observations fixed, comparing \((1000,1000)\), \((400,1600)\)
and \((1600,400)\). Their exact settings are in
Appendix~\ref{app:construction}, with results in
Appendix~\ref{app:additional-results}. A separate earlier jackknife
investigation checks whether changing the MI bias correction helps the
sparse-table failures; its 20,000-pair design is retained in
Appendix~\ref{sec:jackknife-difference-check}.

\section{Sampling and confidence intervals}

The experiment contains 366 distinct evaluation settings. Each uses
100,000 independently sampled pairs of multinomial count tables: one
from \(P\) and another independently from \(Q\). All four rules receive
the same pair. A further 100,000 pairs at each of 186 distinct equal-MI
settings fit the moment references and calibration cutoffs. Separate
random streams supply fitting and evaluation samples. Settings shared
between comparisons are sampled once, giving 55.2 million pairs in this
study. The six upper-null sensitivity settings add 100,000 fitting and
100,000 validation pairs each, for another 1.2 million pairs.

The primary rejection rate divides valid rejections by all 100,000
attempted pairs. Invalid outputs are recorded separately. The bias of
the corrected MI difference is its mean estimate minus the true
difference. Standard-error accuracy is measured by dividing the mean
squared standard error by the observed variance of the difference:
a ratio of one indicates agreement. Under alternatives, the experiment
also records whether rejections estimate the correct direction of the
MI difference. Appendix~\ref{sec:estimation-diagnostics} specifies the
calculations and inclusion rules for all summaries.

Wilson 95\% confidence intervals describe uncertainty in rejection rates
\citep{wilson1927}. Intervals for differences between methods account
for their use of the same sampled pairs. Calibrated-power intervals also
include uncertainty from fitting the cutoffs. A difference within one
percentage point is the stated practical tolerance for the focused
power comparison. The intervals are pointwise: each concerns one rate
or comparison at a specified population setting. Appendix~\ref{sec:simulation-uncertainty} gives the
calculations, including moment uncertainty and the calibration bootstrap.
